\section{Threats to Validity}
\label{sec:threats}

The main threat is overinterpreting the fitted abstraction. The chain
remains a reliability model for specified trace features under
specified diagnostics, not a universal model of agent behavior, so the
estimates should be read conditionally.

\paragraph{Construct validity (state construction).}
Mapping trajectories to a DTMC aggregates over memory, tool state, and
context, and no single mapping is canonical.
\S\ref{sec:state_construction} gives a reproducible construction whose
output is \emph{tested} by Algorithm~\ref{alg:gof}: SS7 checks the
protocol on known synthetic chains and SS9 checks controlled fit/test
recovery. These establish statistical adequacy for the chosen features,
not uniqueness of $\phi$. If GoF rejects, first-passage interpretation
should stop until traces are segmented, re-featurized, or modeled with
a richer process; semi-Markov or hidden Markov models (HMMs) are the
natural next steps. We therefore recommend reporting
$(\phi, m, p_{\mathrm{KS}}, \Delta_{\rm AIC})$ with each estimate.

\paragraph{Estimator sensitivity.}
The composite verdict is invariant across the $18$ configurations of
SS11, and $\hat R_\infty$ varies only over $[0.660,0.674]$ there. Even
so, $R_\infty$ is a property of the split that produced it---sub-sampling
moves the held-out target as well as the estimate---so it must always be
reported against the held-out rate of its own split.

\paragraph{Empirical scope of SS9.}
SS9 uses synthetic traces from a true absorbing chain with noisy
one-hot emissions, so its held-out KS test evaluates controlled
recovery by the trace-to-chain pipeline (featurizer + clustering + MLE)
rather than the full modeling class on unprocessed benchmark data;
\S\ref{sec:ss10}--\ref{sec:ss15} supply the real-trace counterpart.

\paragraph{Independence assumption.}
The pass$^k$ and pass$@k$ identities require i.i.d.\ trials
(A\ref{A:iid}), which can fail under temperature-0 sampling or shared
prefixes. Theorem~\ref{thm:corr} and Corollary~\ref{cor:diag} quantify
the resulting bias and provide a diagnostic.

\paragraph{External validity.}
The MAST case study covers 7 frameworks, so its conclusions may not
generalize to other agents, tasks, or trace distributions; SS6 is
therefore framed as illustrative rather than confirmatory.

\paragraph{NHPP-limit scope.}
The NHPP limit applies under rare-failure scaling
($\varepsilon\to 0$, $d\to\infty$); in high-failure regimes
($\varepsilon > 0.1$) the exact closed form of
Proposition~\ref{thm:closed} should be used instead. That is the common
case for deployed agents---the SWE-agent corpus of \S\ref{sec:ss10}
fails on roughly $32\%$ of tasks ($155/479$)---so the
Goel--Okumoto/NHPP limit is best treated as a conceptual bridge to
reliability-growth theory rather than an estimator at these rates.

%% ================================================================
